% GENERATED FILE. Edit canonical lessons, then run npm run generate:books.
% canonical-source-hash: fnv1a64:b2647362ca8efb95
% canonical-lessons: ML-C349-saksi, ML-C349-yuddham, ML-C349-lekhakan, ML-C349-sambhavikkuka, ML-C349-prakhyapikkuka

\chapter{A Witness to the War}
\label{ch:ml-a-witness-to-the-war}
\hlchaptermodality{349}

\begin{chapteropening}
\textbf{By the end of this chapter:} \emph{I can say a witness, a war, a correspondent, to happen and to announce.}

Complete the last lesson of chapter 349: I can say a witness, a war, a correspondent, to happen and to announce.
\end{chapteropening}

\section[\texorpdfstring{\ml{സാക്ഷി} (sākṣi)}{സാക്ഷി (sākṣi)}]{\ml{സാക്ഷി} (sākṣi) — a witness}

\label{lesson:ML-C349-saksi}



\begin{quote}
Before the new one: say the Malayalam for a principal, then the Malayalam for tuition.
\end{quote}

\subsection*{You'll want to know: \ml{സാക്ഷി}}
\textbf{\ml{സാക്ഷി}} — \emph{sākṣi} — \textquotedblleft{}a witness\textquotedblright{}.

\begin{quote}
\ml{അപകടം} \ml{കണ്ട} \ml{സാക്ഷി} \ml{പോലീസിനോട്} \ml{സംസാരിച്ചു}.
\end{quote}

\begin{quote}
\emph{apakaṭaṁ kaṇṭa sākṣi pōlīsinōṭŭ saṁsāriccu.}
\end{quote}

\begin{quote}
The witness who saw the accident spoke to the police.
\end{quote}

Plural: \ml{സാക്ഷികൾ}. To give evidence is \ml{സാക്ഷി} \ml{പറയുക}, literally to say as a witness.

\begin{cousinweb}
\textbf{\ml{സാക്ഷി}} (\emph{sākṣi}) is borrowed from Sanskrit \emph{sākṣin}, \textquotedblleft{}a witness\textquotedblright{} (CDIAL 13321). CDIAL takes the Sanskrit word apart as \emph{sa-} \textquotedblleft{}with\textquotedblright{} and \emph{akṣi} \textquotedblleft{}eye\textquotedblright{}: the witness is the one who was there with his own eyes.
\end{cousinweb}

\begin{grammarlens}[title={What you've built}]
One more word for telling what happened.
\end{grammarlens}

\subsection*{Your turn}
\begin{itemize}
\raggedright
  \item \emph{Say it:} \emph{sākṣi}
  \item \emph{Say it:} \emph{sākṣi}, once more
  \item \emph{Say it:} \emph{apakaṭaṁ kaṇṭa sākṣi pōlīsinōṭŭ saṁsāriccu.} — The witness who saw the accident spoke to the police.
  \item \emph{Recall:} say the Malayalam for a principal, then the Malayalam for tuition, then say \emph{sākṣi} again
\end{itemize}

\begin{tcolorbox}[breakable,colback=teal!4,colframe=teal!35!black,title={Before you move on}]
What does \ml{സാക്ഷി} mean? (\textquotedblleft{}A witness\textquotedblright{}.) Say it once more.
\end{tcolorbox}
\section[\texorpdfstring{\ml{യുദ്ധം} (yuddhaṁ)}{യുദ്ധം (yuddhaṁ)}]{\ml{യുദ്ധം} (yuddhaṁ) — a war; a battle}

\label{lesson:ML-C349-yuddham}



\begin{quote}
Before the new one: say the Malayalam for tuition, then the Malayalam for a witness.
\end{quote}

\subsection*{You'll want to know: \ml{യുദ്ധം}}
\textbf{\ml{യുദ്ധം}} — \emph{yuddhaṁ} — \textquotedblleft{}a war; a battle\textquotedblright{}.

\begin{quote}
\ml{ആ} \ml{യുദ്ധം} \ml{വളരെ} \ml{കാലം} \ml{മുൻപ്} \ml{നടന്നു}.
\end{quote}

\begin{quote}
\emph{ā yuddhaṁ vaḷare kālaṁ munpŭ naṭannu.}
\end{quote}

\begin{quote}
That war took place a very long time ago.
\end{quote}

With \ml{ചെയ്യുക} it makes a verb: \ml{യുദ്ധം} \ml{ചെയ്യുക}, to fight a war.

\begin{cousinweb}
\textbf{\ml{യുദ്ധം}} (\emph{yuddhaṁ}) is borrowed from Sanskrit \emph{yuddha}, \textquotedblleft{}fought\textquotedblright{}, and as a noun \textquotedblleft{}a battle\textquotedblright{}, the past participle of the root \emph{yudh}, \textquotedblleft{}to fight\textquotedblright{} (CDIAL 10499). Like many Sanskrit nouns in -a, it ends in -\ml{ം} in Malayalam.
\end{cousinweb}

\begin{grammarlens}[title={What you've built}]
One more word for telling what happened.
\end{grammarlens}

\subsection*{Your turn}
\begin{itemize}
\raggedright
  \item \emph{Say it:} \emph{yuddhaṁ}
  \item \emph{Say it:} \emph{yuddhaṁ}, once more
  \item \emph{Say it:} \emph{ā yuddhaṁ vaḷare kālaṁ munpŭ naṭannu.} — That war took place a very long time ago.
  \item \emph{Recall:} say the Malayalam for tuition, then the Malayalam for a witness, then say \emph{yuddhaṁ} again
\end{itemize}

\begin{tcolorbox}[breakable,colback=teal!4,colframe=teal!35!black,title={Before you move on}]
What does \ml{യുദ്ധം} mean? (\textquotedblleft{}A war; a battle\textquotedblright{}.) Say it once more.
\end{tcolorbox}
\section[\texorpdfstring{\ml{ലേഖകൻ} (lēkhakan)}{ലേഖകൻ (lēkhakan)}]{\ml{ലേഖകൻ} (lēkhakan) — a correspondent; a writer}

\label{lesson:ML-C349-lekhakan}



\begin{quote}
Before the new one: say the Malayalam for a witness, then the Malayalam for a war.
\end{quote}

\subsection*{You'll want to know: \ml{ലേഖകൻ}}
\textbf{\ml{ലേഖകൻ}} — \emph{lēkhakan} — \textquotedblleft{}a correspondent; a writer\textquotedblright{}.

\begin{quote}
\ml{പത്രത്തിന്റെ} \ml{ലേഖകൻ} \ml{ഗ്രാമത്തിൽ} \ml{എത്തി}.
\end{quote}

\begin{quote}
\emph{patrattinṟe lēkhakan grāmattil etti.}
\end{quote}

\begin{quote}
The newspaper's correspondent reached the village.
\end{quote}

A woman writer or correspondent is \ml{ലേഖിക}. A correspondent writes about events: \ml{ലേഖകൻ} \ml{ആഘോഷം} \ml{കണ്ടു}, the correspondent saw the celebration.

\begin{cousinweb}
\textbf{\ml{ലേഖകൻ}} (\emph{lēkhakan}) is borrowed from Sanskrit \emph{lekhaka}, \textquotedblleft{}a writer, a scribe, a clerk\textquotedblright{} (Monier-Williams), with the Malayalam masculine ending -\ml{ൻ}. It belongs to the same family as \ml{ലേഖനം} (\emph{lēkhanaṁ}), \textquotedblleft{}an article\textquotedblright{}.
\end{cousinweb}

\begin{grammarlens}[title={What you've built}]
One more word for telling what happened.
\end{grammarlens}

\subsection*{Your turn}
\begin{itemize}
\raggedright
  \item \emph{Say it:} \emph{lēkhakan}
  \item \emph{Say it:} \emph{lēkhakan}, once more
  \item \emph{Say it:} \emph{patrattinṟe lēkhakan grāmattil etti.} — The newspaper's correspondent reached the village.
  \item \emph{Recall:} say the Malayalam for a witness, then the Malayalam for a war, then say \emph{lēkhakan} again
\end{itemize}

\begin{tcolorbox}[breakable,colback=teal!4,colframe=teal!35!black,title={Before you move on}]
What does \ml{ലേഖകൻ} mean? (\textquotedblleft{}A correspondent; a writer\textquotedblright{}.) Say it once more.
\end{tcolorbox}
\section[\texorpdfstring{\ml{സംഭവിക്കുക} (sambhavikkuka)}{സംഭവിക്കുക (sambhavikkuka)}]{\ml{സംഭവിക്കുക} (sambhavikkuka) — to happen; to occur}

\label{lesson:ML-C349-sambhavikkuka}



\begin{quote}
Before the new one: say the Malayalam for a war, then the Malayalam for a correspondent.
\end{quote}

\subsection*{You'll want to know: \ml{സംഭവിക്കുക}}
\textbf{\ml{സംഭവിക്കുക}} — \emph{sambhavikkuka} — \textquotedblleft{}to happen; to occur\textquotedblright{}.

\begin{quote}
\ml{ഇന്നലെ} \ml{രാത്രി} \ml{എന്താണ്} \ml{സംഭവിച്ചത്}?
\end{quote}

\begin{quote}
\emph{innale rātri entāṇŭ sambhaviccatŭ?}
\end{quote}

\begin{quote}
What happened last night?
\end{quote}

Past: \ml{സംഭവിച്ചു}, \emph{sambhaviccu}.

The person something happens to takes -\ml{ക്ക്} or -\ml{ന്}: \ml{അവന്} \ml{എന്ത്} \ml{സംഭവിച്ചു}? What happened to him?

\begin{cousinweb}
\textbf{\ml{സംഭവിക്കുക}} (\emph{sambhavikkuka}) is built from Sanskrit \emph{saṁbhava}, \textquotedblleft{}origin, coming into being\textquotedblright{} (CDIAL 12963), with the Malayalam verb ending -\ml{ിക്കുക}. \ml{സംഭവം} (\emph{sambhavaṁ}), \textquotedblleft{}an event\textquotedblright{}, is the noun of the same word: what happens is a \ml{സംഭവം}.
\end{cousinweb}

\begin{grammarlens}[title={What you've built}]
One more word for telling what happened.
\end{grammarlens}

\subsection*{Your turn}
\begin{itemize}
\raggedright
  \item \emph{Say it:} \emph{sambhavikkuka}
  \item \emph{Say it:} \emph{sambhavikkuka}, once more
  \item \emph{Say it:} \emph{innale rātri entāṇŭ sambhaviccatŭ?} — What happened last night?
  \item \emph{Recall:} say the Malayalam for a war, then the Malayalam for a correspondent, then say \emph{sambhavikkuka} again
\end{itemize}

\begin{tcolorbox}[breakable,colback=teal!4,colframe=teal!35!black,title={Before you move on}]
What does \ml{സംഭവിക്കുക} mean? (\textquotedblleft{}To happen; to occur\textquotedblright{}.) Say it once more.
\end{tcolorbox}
\section[\texorpdfstring{\ml{പ്രഖ്യാപിക്കുക} (prakhyāpikkuka)}{പ്രഖ്യാപിക്കുക (prakhyāpikkuka)}]{\ml{പ്രഖ്യാപിക്കുക} (prakhyāpikkuka) — to announce; to declare}

\label{lesson:ML-C349-prakhyapikkuka}



\begin{quote}
Before the new one: say the Malayalam for a correspondent, then the Malayalam for to happen.
\end{quote}

\subsection*{You'll want to know: \ml{പ്രഖ്യാപിക്കുക}}
\textbf{\ml{പ്രഖ്യാപിക്കുക}} — \emph{prakhyāpikkuka} — \textquotedblleft{}to announce; to declare\textquotedblright{}.

\begin{quote}
\ml{പ്രിൻസിപ്പൽ} \ml{നാളെ} \ml{അവധി} \ml{പ്രഖ്യാപിച്ചു}.
\end{quote}

\begin{quote}
\emph{prinsippal nāḷe avadhi prakhyāpiccu.}
\end{quote}

\begin{quote}
The principal announced a holiday for tomorrow.
\end{quote}

Past: \ml{പ്രഖ്യാപിച്ചു}, \emph{prakhyāpiccu}.

A formal word, common in the news. The noun is \ml{പ്രഖ്യാപനം}, an announcement. \ml{ആഘോഷത്തിന്} \ml{അവധി} \ml{പ്രഖ്യാപിച്ചു}: a holiday was declared for the celebration.

\begin{cousinweb}
\textbf{\ml{പ്രഖ്യാപിക്കുക}} (\emph{prakhyāpikkuka}) is built from Sanskrit \emph{prakhyāpana}, \textquotedblleft{}making known, a report\textquotedblright{} (Monier-Williams), with the Malayalam verb ending -\ml{ിക്കുക}. The Sanskrit noun comes from the verb \emph{pra-khyā}, \textquotedblleft{}to announce, to proclaim\textquotedblright{}.
\end{cousinweb}

\begin{grammarlens}[title={What you've built}]
One more word for telling what happened.
\end{grammarlens}

\subsection*{Your turn}
\begin{itemize}
\raggedright
  \item \emph{Say it:} \emph{prakhyāpikkuka}
  \item \emph{Say it:} \emph{prakhyāpikkuka}, once more
  \item \emph{Say it:} \emph{prinsippal nāḷe avadhi prakhyāpiccu.} — The principal announced a holiday for tomorrow.
  \item \emph{Recall:} say the Malayalam for a correspondent, then the Malayalam for to happen, then say \emph{prakhyāpikkuka} again
\end{itemize}

\begin{tcolorbox}[breakable,colback=teal!4,colframe=teal!35!black,title={Before you move on}]
What does \ml{പ്രഖ്യാപിക്കുക} mean? (\textquotedblleft{}To announce; to declare\textquotedblright{}.) Say it once more.
\end{tcolorbox}
